\documentclass[11pt,a4paper]{article}
\usepackage[margin=0.8in]{geometry}
\usepackage{enumitem}
\usepackage{listings}
\usepackage{xcolor}
\usepackage{booktabs}
\usepackage{longtable}
\usepackage{array}
\usepackage{fancyhdr}
\usepackage{titlesec}

\setlength{\headheight}{14pt}
\addtolength{\topmargin}{-2pt}

\pagestyle{fancy}
\fancyhf{}
\fancyhead[L]{Class XI - Computer Science (083)}
\fancyhead[R]{Topic: Tuple in Python}
\fancyfoot[C]{\thepage}
\renewcommand{\headrulewidth}{0.4pt}

\titleformat{\section}{\large\bfseries}{Q\thesection.}{0.5em}{}

\lstset{
    language=Python,
    basicstyle=\ttfamily\footnotesize,
    breaklines=true,
    frame=single,
    backgroundcolor=\color{gray!5},
    keywordstyle=\color{blue}\bfseries,
    stringstyle=\color{red},
    commentstyle=\color{green!60!black},
    showstringspaces=false,
    tabsize=4,
    numbers=left,
    numberstyle=\tiny\color{gray},
    numbersep=3pt,
    framexleftmargin=1.5em
}

\setlist[enumerate]{leftmargin=*,itemsep=0.3em,topsep=0.2em}
\setlist[enumerate,2]{label=\alph*.,leftmargin=1.5em}

\title{\textbf{Tuple Worksheet Questions}\\[0.3em]
\large Class XI Computer Science (083) | Topic: Tuples in Python}
\author{}
\date{}

\begin{document}
\maketitle
\thispagestyle{fancy}

\begin{center}
\textbf{Instructions:} Write the answers in the space provided or on a separate sheet as directed by your teacher.
\end{center}

\vspace{0.5em}

\begin{enumerate}[label=\textbf{Q\arabic*.}]

% Q1
\item What do you mean by Tuple in Python?

% Q2
\item Write a statement to create an empty tuple named \texttt{T1}.

% Q3
\item Write a statement to create a tuple \texttt{T1} containing first five even numbers.

% Q4
\item Write the code to convert the given list \texttt{L1} to tuple.\\
\texttt{L1 = [1, 2, 3, 4, 5]}

% Q5
\item Write a statement to create tuple \texttt{T1} with single element.

% Q6
\item Write the output of the following code:
\begin{lstlisting}
T1 = (4)
print(type(T1))
\end{lstlisting}

% Q7
\item Write the output of the following code:
\begin{lstlisting}
T1 = 4, 5, 6
print(type(T1))
\end{lstlisting}

% Q8
\item Write the output of the following code:
\begin{lstlisting}
T1 = (1, 2, 3, 4, 5, 6, 7, 8)
print(T1)
print(T1 * 2)
print(T1 + T1)
print(len(T1) * 2)
\end{lstlisting}

% Q9
\item Write a statement to create tuple \texttt{T1} with the following data:\\
\texttt{1, 4, 'CS', 'IP', 5}

% Q10
\item Write the output of the following code:
\begin{lstlisting}
T1 = (1, 2, 3, 4, 5, 6, 7, 8)
print(T1[0])
print(T1[-1])
print(T1[2 + 3])
print(T1[4 - 1])
print(T1[7 % 2])
\end{lstlisting}

% Q11
\item Write the output of the following:
\begin{lstlisting}
T1 = (1, 2, 3, 4, 5, 6, 7, 8)
print(T1[:])
print(T1[3:])
print(T1[:4])
print(T1[-2:-5])
\end{lstlisting}

% Q12
\item Write the output of the following:
\begin{lstlisting}
T1 = (1, 2, 3, 4, 5, 6, 7, 8)
print(T1[1::2])
print(T1[-1:-5:-2])
print(T1[::-1])
print(T1[:7:2])
\end{lstlisting}

% Q13
\item Which error is returned by the following code:
\begin{lstlisting}
T1 = (1, 2, 3, 4, 5, 6, 7, 8)
print(T1[15])
\end{lstlisting}

% Q14
\item Write two differences between Tuple and List.

% Q15
\item Write a statement in Python to concatenate the given tuples.\\
\texttt{T1 = (1, 2, 3)}\\
\texttt{T2 = (5, 6, 7)}

% Q16
\item Write the output of the following code:
\begin{lstlisting}
T1 = (45, 67, 98)
T1 = T1 + (1, 2, 3)
print(T1)
\end{lstlisting}

% Q17
\item Write the output of the following code:
\begin{lstlisting}
T1 = (45, 67, 98)
T1 = T1 * 3
print(T1)
\end{lstlisting}

% Q18
\item Write the output of the following code:
\begin{lstlisting}
T1 = (45, 67, 98)
T2 = (45, 67, 98)
print(T1 in T2)
print(45 in T2)
print(45 in T1)
print(T1 + T2)
\end{lstlisting}

% Q19
\item Explain the following functions in reference to Tuple in Python:
\begin{enumerate}
    \item \texttt{len()}
    \item \texttt{count()}
\end{enumerate}

% Q20
\item Write a program to accept five numbers from the user and store it in a tuple \texttt{'T1'}.

% Q21
\item Write a program to accept five fruit names from the user and store it in a tuple \texttt{'F1'}.

% Q22
\item Write a program to store \texttt{'n'} number of sports in a tuple \texttt{'S1'}. (Accept \texttt{'n'} from the user)

% Q23
\item Consider the following tuple and write the code for the following statements:\\
\texttt{T1 = (12, 3, 45, 'Hockey', 'Anil', ('a', 'b'))}
\begin{enumerate}
    \item Display the first element of \texttt{'T1'}
    \item Display the last element of \texttt{'T1'}
    \item Display \texttt{'T1'} in reverse order.
    \item Display \texttt{'Anil'} from tuple \texttt{'T1'}
    \item Display \texttt{'b'} from tuple \texttt{'T1'}
\end{enumerate}

% Q24
\item Write the output of the following:
\begin{lstlisting}
for i in tuple("SHIMLA"):
    print(i + i)
\end{lstlisting}

% Q25
\item Write the output of the following:
\begin{lstlisting}
>>> 7 in tuple('123456789')
\end{lstlisting}

% Q26
\item Write the output of the following:
\begin{lstlisting}
T1 = (23, 32, 4, 5, 2, 12, 23, 7, 9, 10, 23)
print(len(T1), *T1[-1])
print(T1[T1.count(23) + len(T1) - 5])
print(T1.count(T1[6]))
print(T1.count(T1[-1]))
\end{lstlisting}

% Q27
\item Write the output of the following:
\begin{lstlisting}
T1 = (23, 32, 4, 5, 2, 12, 23, 7, 9, 10, 23)
print(max(T1))
print(min(T1))
\end{lstlisting}

% Q28
\item Write the output of the following:
\begin{lstlisting}
T1 = ('Hockey', 'Cricket', 'Football')
print(max(T1))
print(min(T1))
\end{lstlisting}

% Q29
\item Write the output of the following:
\begin{lstlisting}
t1 = ('Raman', 'Ram', 'Ramaiya')
print(max(t1))
print(min(t1))
\end{lstlisting}

% Q30
\item Write the output of the following:
\begin{lstlisting}
T1 = (23, 32, 4, 5, 2, 12, 23, 7, 9, 10, 23)
print(sorted(T1))
print(sorted(T1[2:7]))
print(T1.index(23))
print(T1.index(23, 3, 9))
\end{lstlisting}

% Q31
\item Write the output of the following:
\begin{lstlisting}
T1 = (23, 32, 4, 5, 2, 12, 23, 7, 9, 10, 23)
print(sum(T1) + min(T1))
\end{lstlisting}
\textit{(Note: Original may have had a typo in map(); corrected for sense)}

% Q32
\item Write a program to accept three numbers from the user and insert it at the end of given Tuple \texttt{T1}.\\
\texttt{T1 = (23, 32, 4, 5, 2, 12, 23, 7, 9, 10, 23)}

% Q33
\item What type of error is returned by following statement:
\begin{lstlisting}
>>> (a, b, c, d) = (5, 6, 8)
\end{lstlisting}

% Q34
\item Write the output of the following code:
\begin{lstlisting}
>>> (name, rollNo, subject) = ('Anil', 9, 'CS')
>>> name
\end{lstlisting}

% Q35
\item Predict the output:
\begin{lstlisting}
T = (1, 2, 3)
del T
print(T)
\end{lstlisting}

% Q36
\item What do you mean by Nested Tuple? Give one example.

% Q37
\item Write a statement to print 30 from the given tuple.\\
\texttt{a = ("Seventy", [1, 2, 3], (20, 30, 40), "Eighty")}

% Q38
\item Write a statement to unpack the following tuple into 3 variables.\\
\texttt{(60, 70, 80)}

% Q39
\item Write a program to swap the values of given tuples.\\
\texttt{T1 = ('A', 'B')}\\
\texttt{T2 = (34, 65)}

% Q40
\item Write the output of the following code:
\begin{lstlisting}
T1 = ('A', ['B', 'D'], 'C')
T1[1][1] = 'C'
print(T1)
\end{lstlisting}

% Q41
\item Write a program to print the frequency of a number accepted from the user in given tuple:\\
\texttt{T1 = (12, 17, 18, 25, 19, 12, 18, 5)}

% Q42
\item Write a program in Python to concatenate all the characters of given tuple.\\
\texttt{T1 = ('B', 'O', 'O', 'K')}\\
\textit{Expected OUTPUT : BOOK}

% Q43
\item Write a program to remove a number (accepted from the user) from the given tuple\\
\texttt{T1 = (12, 15, 18, 21, 24, 27, 30)}\\
\textit{SAMPLE EXECUTION}\\
\textit{Enter element to remove : 21}\\
\textit{Tuple after removing element is : (12, 15, 18, 24, 27, 30)}

% Q44
\item What do you mean by Unpacking Tuple? Give example.

% Q45
\item What do you mean by Packing Tuple? Give Example.

% Q46
\item Write one similarity between Tuples and String.

% Q47
\item Write the output of the following:
\begin{lstlisting}
T1 = (2, 'Apple', 6)
print(max(T1))
\end{lstlisting}

% Q48
\item Consider the given tuple and write the output of given statements:\\
\texttt{T1 = (1, 23, 4, 5, 'A', ['C', 'D'], (23, 45), 45)}
\begin{enumerate}
    \item \texttt{print(len(T1))}
    \item \texttt{print(T1.index(45))}
    \item \texttt{print(T1.count(45))}
    \item \texttt{print(T1[5][0] - 2)}   % Note: list index, but as per question
    \item \texttt{print(T1[:5])}
    \item \texttt{print(T1[5:])}
    \item \texttt{print(T1.index("A"))}
    \item \texttt{print(T1[-1:-7:-2])}
    \item \texttt{print(T1[4] + T1[5])}
    \item \texttt{print(max(T1))}
\end{enumerate}

% Q49
\item Write the length of following tuple:\\
\texttt{T1 = (1, (2,3), ((34,56),45,67), 39)}

% Q50
\item Write a program to store the detail (Roll number, Name, Section) of three students in a tuple. Accept all data from the user.\\
\textit{SAMPLE OUTPUT :}\\
\texttt{((1, 'Amita', 'B'), (2, 'Sunita', 'C'), (3, 'Daya', 'A'))}

% Q51
\item Write a program to accept roll number from user and display it's detail from given tuple\\
\texttt{Det = ((1, 'Amita', 'B'), (2, 'Sunita', 'C'), (3, 'Daya', 'A'))}

% Q52
\item Write a program to accept five numbers from the user and store these numbers in a tuple. Display the following from tuple:
\begin{enumerate}
    \item Largest number
    \item Smallest number
    \item Sum of all numbers
    \item Average of all numbers
\end{enumerate}

% Q53
\item Write a program in Python to print the second largest element in tuple given below:\\
\texttt{T1 = (23, 45, 87, 45, 67, 43, 23, 12)}

% Q54
\item Write a program in Python to display the sum of all even numbers and odd numbers separately from the tuple given below:\\
\texttt{T1 = (24, 45, 87, 46, 67, 44, 23, 12)}

% Q55
\item Write the output of the following:
\begin{lstlisting}
T1 = ("Amit", "Sumit")
T2 = "Sumit", "Amit"
T3 = "Amit", "Sumit"
print(T1 == T2)
print(T1 == T3)
\end{lstlisting}

% Q56
\item Write a program to accept a number from the user and store first 10 multiples of number in a tuple.\\
\textit{For example:}\\
\textit{Enter any number : 4}\\
\textit{Tuple = (4, 8, 12, 16, 20, 24, 28, 32, 36, 40)}

% Q57
\item Write a program to accept a string from the user and convert it into tuple.\\
\textit{For example:}\\
\textit{Enter any String : tuple}\\
\textit{Tuple = ('t', 'u', 'p', 'l', 'e')}

% Q58
\item Write a program to display the sum of all the numbers in a given tuple:\\
\texttt{T1 = (1, 2, 3, 'A', 'B', 4, 'Sun')}

% Q59
\item Write a program to find the sum of marks of all students stored in the given tuple.\\
\texttt{T1 = (("Suman", 75), ("Glory", 35), ("Ravi", 50))}\\
\textit{Expected Output is : 160}

% Q60
\item Write the output of the following:
\begin{lstlisting}
T1 = (1, (2, 'a'), (3, 'b', 'c'), (7, 'd', 'e', 'f'))
print(len(T1))
print(T1[1][0])
print(7 in T1[3])
print(T1.index(1))
print(len(T1[3]))
print(T1[T1[2][0]])
\end{lstlisting}

% Q61
\item Write the output of the following:
\begin{lstlisting}
T1 = (1, (2, 'a', (3, 'b', 'c', (7, 'd', 'e', 'f'))))
print(len(T1))
print(T1[1][0])
print(7 in T1[1])
print(T1.index(1))
print(len(T1[1]))
print(T1[T1[0]])
\end{lstlisting}

% Q62
\item Write the output of the following:
\begin{lstlisting}
T1 = (1, (2, 'a') * 2)
print(T1)
print(len(T1))
print(T1[1][1])
print(len(T1[1]))
\end{lstlisting}

% Q63
\item Write the output of the following:
\begin{lstlisting}
T1 = (1, (2, 'a')) * 2
print(T1)
print(len(T1))
print(T1[1][1])
print(len(T1[1]))
print(T1.count(1))
print(T1.index(1))
\end{lstlisting}

% Q64
\item Write the output of the following:
\begin{lstlisting}
T1 = ('a') * 3
T2 = ('a',) * 3
print(T1)
print(type(T1))
print(len(T1))
print(T2)
print(type(T2))
print(len(T2))
\end{lstlisting}

% Q65
\item Write the output of the following:
\begin{lstlisting}
T1 = "Spider Man:No Way Home"
T2 = tuple(T1.split())
print(T1)
print(type(T1))
print(len(T1))
print(T2)
print(T2[0][0] + T2[1][0])
print(len(T2))
\end{lstlisting}

% Q66
\item What type of error is returned by following code:
\begin{lstlisting}
T1 = (1, 2, 3, 4)
T2 = (4, 5, 6)
print(T1 - T2)
\end{lstlisting}

% Q67
\item Write the output of the following:
\begin{lstlisting}
T1 = (1, 2, 3, 4)
print(T1 * 3)
print(T1 * (3))
print((T1) * 3)
print(T1 + (3,))
\end{lstlisting}

% Q68
\item Write the output of the following:
\begin{lstlisting}
T1 = (30, 20, 30)
print(T1, T1, T1)
print(T1 * 3)
\end{lstlisting}

% Q69
\item Write the output of the following:
\begin{lstlisting}
a = 30
b = (30)
c = (30,)
print(b)
print(type(b))
print(c)
print(type(c))
\end{lstlisting}

% Q70
\item Write a program to accept a number from the user and create a tuple containing all factors of the number.\\
\textit{Sample Output:}\\
\textit{Enter any Number : 9}\\
\textit{Tuple is : (1, 3, 9)}

% Q71
\item Write a program to arrange all the sub-tuples of a tuple in increasing order and store them in another tuple.\\
\textit{For example:}\\
\texttt{T1 = ((1, 2, 3), (45, 23), (98, 34, 67))} \textit{\# Original Tuple}\\
\texttt{T2 = ((1, 2, 3), (23, 45), (34, 67, 98))} \textit{\# New Tuple}

% Q72
\item Write a program to store the length of all sub-tuples of a tuple into a new tuple.\\
\textit{For example:}\\
\texttt{T1 = ((1, 2, 3), (45, 23), (98, 34, 67), (34, 56, 78, 9))} \textit{\# Original}\\
\texttt{T2 = (3, 2, 3, 4)} \textit{\# New Tuple}

% Q73
\item Write a program to print square of all odd numbers of a tuple.\\
\texttt{T1 = (12, 13, 22, 7, 9, 66, 4)}\\
\textit{Expected Output : 169, 49, 81}

% Q74
\item Write a program to print the common elements of given tuples:\\
\texttt{T1 = (1, 2, 3, 4, 5, 6)}\\
\texttt{T2 = (3, 5, 7, 9, 11)}\\
\textit{Expected Output : 3, 5}

% Q75
\item Write a program to find the union (common elements included once) of given tuples:\\
\texttt{t1 = (1, 2, 3, 4, 5, 6)}\\
\texttt{t2 = (3, 5, 7, 9, 11)}\\
\textit{Expected Output : 1, 2, 3, 4, 5, 6, 7, 9, 11}

% Q76
\item Write a program to store the sum of all sub-tuples of given tuple into a new tuple.\\
\textit{For example:}\\
\texttt{T1 = ((1, 2, 3), (45, 23), (98, 34, 67), (34, 56, 78, 9))} \textit{\#Original}\\
\texttt{T2 = (6, 68, 199, 177)} \textit{\# New Tuple}

% Q77
\item Write a program to store the maximum value of all sub-tuples of given tuple into a new tuple.\\
\textit{For example:}\\
\texttt{T1 = ((1, 2, 3), (45, 23), (98, 34, 67), (34, 56, 78, 9), ("A", "B", "c"))}\\
\texttt{T2 = (3, 45, 98, 78, 'c')} \textit{\# New Tuple (note: string max)}

% Q78
\item Write a program to display the longest sub-tuple from the given tuple.\\
\texttt{T1 = ((1, 2, 3), (45, 23), (98, 34, 67), (34, 56, 78, 9))}\\
\texttt{T2 = (34, 56, 78, 9)} \textit{\# New Tuple}

% Q79
\item Write a program to display names starting from vowel from the given tuple.\\
\texttt{T1 = ("Amit", "Ram", "Esha", "Harry")}

% Q80
\item Write a program to print the length of all names in a given tuple.\\
\texttt{T1 = ("Amit", "Ram", "Esha", "Harry")}\\
\textit{Expected Output : 4, 3, 4, 5}

% Q81
\item What is the purpose of following operators in reference to tuples.
\begin{enumerate}
    \item \texttt{+}
    \item \texttt{*}
    \item \texttt{in}
\end{enumerate}

% Q82
\item Match the following:

\begin{center}
\begin{tabular}{|l|l|l|}
\hline
\textbf{Functions} & \textbf{Description} & \textbf{Match} \\
\hline
\texttt{max()} & Returns the length of the tuple. & (3) \\
\texttt{min()} & Returns the index value of element in tuple. & (4) \\
\texttt{count()} & Returns the largest element in tuple. & (1) \\
\texttt{index()} & Returns the smallest element in tuple. & (2) \\
\texttt{len()} & Returns the frequency of an element in tuple. & (5) \\
\hline
\end{tabular}
\end{center}
\textit{(Match the numbers correctly in your answer)}

% Q83
\item Write the output of the following:
\begin{lstlisting}
>>> S = 1, (2,3,4,5), 51, (61,71)
>>> len(S)
>>> S[2:3]
>>> S[-1][0]
>>> type(S)
\end{lstlisting}

% Q84
\item Predict the output:
\begin{lstlisting}
T = (1, 2, 4, 3, 8, 9)
[T[i] for i in range(0, len(T), 2)]
\end{lstlisting}

% Q85
\item What will be the output of the following code snippet?
\begin{lstlisting}
init_tuple = ()
print(init_tuple.__len__())
\end{lstlisting}

\end{enumerate}

\vspace{1em}
\begin{center}
\textit{--- End of Questions ---}
\end{center}

\end{document}